\chapter{Il sistema a moduli}
\label{cap:08-sistema-moduli}

Il \cref{cap:03-requisiti-principi} ha enunciato la regola che governa la
divisione fra nucleo e moduli: il binario prodotto senza moduli non deve
contenerli. Questo capitolo descrive il meccanismo che la realizza, ne discute i
limiti, e presenta il criterio con cui il progetto tratta le dipendenze fra
moduli.

\section{Perch\'e a tempo di build e non a runtime}

L'approccio convenzionale a un sistema estensibile \`e l'attivazione
condizionale: tutte le funzionalit\`a sono compilate nel binario e un flag di
configurazione decide quali istanziare. Spring Boot lo supporta nativamente con
\code{@ConditionalOnProperty}, ed \`e la scelta che quasi ogni applicazione fa.

\nf{} non la fa, per tre ragioni.

\paragraph{La superficie del binario.} Un modulo disattivato \`e comunque
presente: le sue classi sono nel classpath, le sue dipendenze transitive sono
risolte, il suo codice \`e analizzato dall'ottimizzatore. Per la compilazione
in immagine nativa (\cref{cap:20-build-packaging}) questo \`e sostanziale: il
provider Kubernetes trascina il client Fabric8 e Vert.x, il provider container
trascina il client Docker Java. Un'immagine nativa che contenga entrambi \`e
significativamente pi\`u grande e pi\`u lenta da compilare di una che contenga
solo ci\`o che verr\`a usato.

\paragraph{L'onest\`a della misura.} Un esperimento che confronta il costo del
control plane con e senza controllo di concorrenza deve confrontare due binari
diversi, non lo stesso binario con un flag diverso. Nel secondo caso le classi
inattive contribuiscono comunque all'impronta di memoria, al tempo di
inizializzazione del contesto Spring e alla pressione sulla cache di istruzioni.
Il \cref{cap:24-footprint} misura configurazioni, e questo richiede che le
configurazioni siano artefatti distinti.

\paragraph{L'assenza di rami condizionali.} \`E l'argomento pi\`u forte. Se i
moduli fossero attivabili a runtime, il codice del nucleo dovrebbe contenere test
sulla loro presenza. Con la selezione a build time, il nucleo dipende solo da
interfacce di cui esiste sempre un'implementazione, e il percorso di invocazione
non contiene un solo \code{if} che riguardi la configurazione modulare.

\section{Il meccanismo}

\subsection{Selezione}

La selezione avviene attraverso una propriet\`a Gradle o una variabile d'ambiente:

\begin{lstlisting}[language=bash]
./gradlew :control-plane:bootJar -PcontrolPlaneModules=none
./gradlew :control-plane:bootJar -PcontrolPlaneModules=async-queue,sync-queue
./gradlew :control-plane:bootJar -PcontrolPlaneModules=all
NANOFAAS_CONTROL_PLANE_MODULES=all ./gradlew :control-plane:bootJar
\end{lstlisting}

I moduli disponibili sono scoperti scandendo le sottodirectory di
\code{platform/modules/} che contengono uno script di build --- non esiste un
elenco da mantenere allineato. La selezione \`e poi validata con tre regole:

\begin{lstlisting}[language=Java, caption={Validazione della selezione in
\code{settings.gradle}.}]
if (requestedModules.contains('none')) {
    if (requestedModules.size() > 1) {
        throw new GradleException("Module selector 'none' cannot be combined with other values: ...")
    }
    requestedModules.clear()
}
if (requestedModules.contains('all')) {
    requestedModules = new LinkedHashSet<String>(availableModules)
}
def unknownModules = requestedModules.findAll { !availableModules.contains(it) }
if (!unknownModules.isEmpty()) {
    throw new GradleException("Unknown control-plane modules: ${unknownModules}. Available: ${availableModules}")
}
\end{lstlisting}

Un nome sconosciuto fa fallire la build. \`E la scelta corretta per un selettore
di questo tipo: ignorare silenziosamente \code{-PcontrolPlaneModules=async-quee}
produrrebbe un binario privo della coda asincrona e un operatore convinto del
contrario.

I moduli selezionati entrano nel classpath come dipendenze \code{runtimeOnly}.
La scelta \`e significativa: il codice del control plane non pu\`o compilare
contro una classe di modulo, perch\'e a tempo di compilazione i moduli non sono
sul suo compile classpath. Il compilatore \`e cos\`i l'esecutore della regola di
isolamento, non un revisore.

\subsection{Caricamento}

A runtime, i moduli si presentano attraverso il \code{ServiceLoader} di Java.
Ciascuno pubblica un file
\code{META-INF/services/it.unimib.datai.nanofaas.common.controlplane.ControlPlaneModule}
contenente il nome della propria classe di implementazione, che \`e tipicamente
di dieci righe:

\begin{lstlisting}[language=Java]
public final class ConcurrencyControlModule implements ControlPlaneModule {
    @Override
    public Set<Class<?>> configurationClasses() {
        return Set.of(ConcurrencyControlConfiguration.class);
    }
}
\end{lstlisting}

Il ponte fra \code{ServiceLoader} e Spring \`e un \code{DeferredImportSelector}:

\begin{lstlisting}[language=Java, caption={\code{ControlPlaneModuleImportSelector}.}]
public final class ControlPlaneModuleImportSelector implements DeferredImportSelector {
    @Override
    public String[] selectImports(AnnotationMetadata importingClassMetadata) {
        ClassLoader classLoader = Thread.currentThread().getContextClassLoader();
        Set<String> moduleConfigurations = new LinkedHashSet<>();
        ServiceLoader.load(ControlPlaneModule.class, classLoader)
                .forEach(module -> module.configurationClasses().stream()
                        .filter(Objects::nonNull)
                        .map(Class::getName)
                        .forEach(moduleConfigurations::add));
        return moduleConfigurations.toArray(String[]::new);
    }
}
\end{lstlisting}

Il commento nel sorgente ne spiega la ragione d'essere: \emph{``Ensures optional
module \code{@Configuration} classes are imported for every bootstrap path,
including tests that do not execute \code{ControlPlaneApplication.main()}.''} Un
caricamento eseguito nel \code{main} funzionerebbe in produzione e non nei test
di integrazione, che avviano il contesto per altra via. Legandolo a
un'annotazione \code{@Import} sulla classe applicativa, ogni percorso di
avvio che usi quella classe carica i moduli.

Va notato che i package dei moduli \textbf{non} sono soggetti a component
scanning: ogni modulo registra i propri bean con \code{@Bean} espliciti nella
propria classe di configurazione. La scansione automatica avrebbe reso il
contenuto del contesto dipendente dai nomi dei package, che \`e esattamente il
tipo di accoppiamento implicito che il sistema modulare vuole evitare.

\section{Le dipendenze fra moduli}

L'interfaccia \code{ControlPlaneModule} non consente a un modulo di dichiarare
di averne bisogno di un altro. \`E una lacuna apparente, e il progetto l'ha
mantenuta di proposito: l'unica dipendenza legittima fra moduli \`e attraverso
un'interfaccia del nucleo, e in tal caso il modulo dipende dall'interfaccia,
non dal modulo che la implementa.

Rimane per\`o il caso in cui l'implementazione neutra dell'interfaccia rende il
modulo dipendente \emph{inerte}, e questo caso ha bisogno di una risposta.

\subsection{Il caso fatale: \code{concurrency-control}}

Il modulo \code{concurrency-control} legge lo stato delle code da
\code{ScalingMetricsSource} e vi \emph{riscrive} i limiti che calcola. Il nucleo
registra sempre un'implementazione no-op di quell'interfaccia, quindi la
condizione \code{@ConditionalOnBean} \`e soddisfatta anche quando nessun modulo
ne fornisce una reale. Il costruttore della configurazione codifica quindi il
controllo che l'annotazione non pu\`o fare:

\begin{lstlisting}[language=Java]
/**
 * Refuses to start when no module supplies real queue state.
 *
 * <p>{@code @ConditionalOnBean} above cannot catch this: the core always registers a no-op
 * source, so the condition holds even when nothing produces one. The governor would then run
 * against a source reporting depth 0 and in-flight 0 forever, and -- worse -- every limit it
 * computed would be written into that same no-op and enforced by nobody. The module would be
 * entirely inert while its metrics claimed otherwise.
 */
public ConcurrencyControlConfiguration(ScalingMetricsSource metricsSource) {
    if (!metricsSource.enabled()) {
        throw new IllegalStateException(
                "concurrency-control needs a module that supplies queue state (async-queue); "
                        + "without one the governor reads zeroes and the limits it computes "
                        + "are enforced by nothing. Add async-queue to the module selection.");
    }
}
\end{lstlisting}

Il metodo \code{enabled()} sull'interfaccia \`e ci\`o che rende possibile la
distinzione. Un'implementazione no-op \`e indistinguibile da una reale che
osserva un sistema scarico --- entrambe riportano zero --- a meno che non
dichiari di essere no-op. \`E una tecnica generale: quando un default neutro
pu\`o essere confuso con un'osservazione legittima, deve poter essere
interrogato sulla propria natura.

Il messaggio di errore nomina il modulo mancante e l'azione correttiva. In un
sistema in cui la configurazione avviene a tempo di build, un errore di
configurazione si manifesta all'avvio ed \`e l'unica occasione per dire
all'operatore che cosa fare.

\subsection{Il caso non fatale: \code{autoscaler}}

Lo stesso modulo \code{autoscaler} legge la stessa sorgente, ma non fallisce.
La differenza \`e che la sua metrica principale --- \code{rps} --- proviene da
un contatore Micrometer indipendente dalle code. Senza \code{async-queue}
l'autoscaler resta pienamente funzionante sulla metrica di default, e sono le
due metriche derivate dalle code (\code{queue\_depth} e \code{in\_flight}) a
diventare cieche. Il modulo emette un avviso al primo uso di una di esse.

Il criterio che distingue i due casi \`e quello enunciato nel
\cref{cap:03-requisiti-principi}: degradazione osservabile, s\`i; degradazione
silenziosa, no. \code{autoscaler} degrada parzialmente e lo dichiara;
\code{concurrency-control} degraderebbe totalmente e in silenzio, quindi non
degrada affatto.

\section{Isolamento verificato}

La regola secondo cui un modulo non deve dipendere da un altro modulo \`e
verificata automaticamente con ArchUnit~\cite{archunit}, e ciascuno dei nove
moduli porta la propria classe di verifica. La regola \`e in vigore senza
eccezioni: al momento della stesura nessun modulo dichiara una deroga.

La formulazione della regola merita una nota, perch\'e definisce con precisione
che cosa sia una dipendenza illegittima. Un modulo non pu\`o dipendere da
\emph{altri moduli}; pu\`o invece dipendere liberamente dalle interfacce del
nucleo, che \`e il meccanismo previsto. Cos\`i i provider di deployment
implementano \code{ImageValidator}, dichiarata nel package \code{registry} del
control plane (\cref{cap:15-deployment-provider}), senza violare nulla: il
riferimento va verso il nucleo, non lateralmente.

Il costo di questa formulazione \`e che una dipendenza fra moduli non pu\`o
essere introdotta per errore ma pu\`o esserlo per decisione: il commento nel
codice prescrive che ogni coppia sanzionata sia aggiunta come riga esplicita con
la propria giustificazione. Nessuna \`e stata aggiunta finora.

\section{Il catalogo dei moduli}

\begin{table}[htbp]
\centering
\caption{I nove moduli opzionali, con dimensione del codice di produzione e
capitolo di riferimento.}
\label{tab:moduli}
\small
\begin{tabularx}{\textwidth}{@{}lrXl@{}}
\toprule
\textbf{Modulo} & \textbf{LOC} & \textbf{Ruolo} & \textbf{Cap.} \\
\midrule
\code{async-queue} & 682 & Code per funzione e scheduler del percorso asincrono
& \ref{cap:09-async-queue} \\
\code{sync-queue} & 905 & Ammissione e backpressure sul percorso sincrono
& \ref{cap:10-sync-queue} \\
\code{autoscaler} & 754 & Scaling interno del numero di repliche
& \ref{cap:11-autoscaler} \\
\code{concurrency-control} & 1\,314 & Governo del limite di concorrenza per
funzione & \ref{cap:12-concurrency-control} \\
\code{offload} & 230 & Proxy condizionale verso un'istanza remota
& \ref{cap:13-offload} \\
\code{runtime-config} & 540 & Riconfigurazione a caldo e API di
amministrazione & \ref{cap:14-runtime-config} \\
\code{k8s-deployment-provider} & 1\,202 & Backend di deployment Kubernetes
& \ref{cap:15-deployment-provider} \\
\code{container-deployment-provider} & 1\,166 & Backend di deployment su runtime
container locale & \ref{cap:15-deployment-provider} \\
\code{build-metadata} & 43 & Endpoint diagnostico sulla build
& \ref{cap:16-moduli-diagnostici} \\
\midrule
\textbf{Totale moduli} & \textbf{6\,836} & & \\
\textit{Nucleo, per confronto} & \textit{5\,065} & & \\
\bottomrule
\end{tabularx}
\end{table}

Il rapporto fra le due ultime righe \`e l'osservazione con cui vale la pena
chiudere il capitolo. La somma dei moduli supera il nucleo che estendono: la
maggior parte del codice di \nf{} \`e codice che si pu\`o non compilare. Se il
meccanismo di selezione modulare fosse mal progettato, questo sarebbe un
problema di manutenzione; essendo verificato dal compilatore e dalle regole
architetturali, \`e invece la misura di quanto la piattaforma sia riducibile.

\section{Il default \`e \code{all}}

Un'ultima nota, di metodo pi\`u che di architettura. In assenza di selezione
esplicita, ogni invocazione di Gradle --- compresa quella dei test --- usa
\code{all}. La ragione \`e che la suite di test deve esercitare la
configurazione che viene distribuita, non una configurazione parziale che
nessuno esegue.

La configurazione priva di moduli non resta per\`o senza copertura: la classe
\code{CoreOnlyApiTest} verifica le propriet\`a che valgono \emph{solo} in
assenza di moduli --- accodatore no-op, \code{501} su \code{:enqueue} --- e si
auto-esclude quando i moduli sono presenti. Poich\'e la selezione modulare \`e
risolta a tempo di configurazione di Gradle e non pu\`o variare all'interno di
una singola build, quel test richiede una propria invocazione con
\code{-PcontrolPlaneModules=none}. \`E un costo reale sulla catena di
integrazione continua, ed \`e il prezzo diretto della scelta di selezionare a
build time anzich\'e a runtime.
